% =====================================================================
%  05-documentos.tex - Verificacion del repositorio documental
% =====================================================================

\section{Verificación del repositorio documental}\label{sec:pruebas-documentos}

Esta sección reúne la verificación del repositorio documental de la \cref{sec:impl-documentos}: las consultas, las versiones inmutables, la clasificación con carga atómica, la admisión general de nuevas cargas, la entrega del contenido exacto con su buzón de integridad y la firma documental individual. Las pruebas nativas y de formato de la política PDF/DOCX de soportes procesales (\cref{sec:impl-formatos}) se presentan con las etapas (\cref{sec:pruebas-etapas}). La trazabilidad por archivo de consultas, versiones, clasificación, contenido y firma se reúne en la \cref{sec:matriz-grupo-documentos}.

\subsection{Consultas documentales e integración de navegador}\label{sec:pruebas-consultas}

Las consultas de metadatos incorporan controles de entrada, autorización, proyección SQL y auditoría. La \cref{tab:pruebas-consultas} resume las condiciones comprobadas por sus pruebas, incluidas en las 539 aprobadas el 14 de septiembre (\cref{tab:pruebas-cobertura-consultas}). El conteo histórico y su cobertura permanecen identificados con la ejecución en que fueron medidos.

{\emergencystretch=2em Los límites de parámetros se prueban en \rutarepo{crates/application/tests/document_query.rs}; la autorización del caso de uso, en \rutarepo{crates/application/tests/case_document_queries.rs}. La integración PostgreSQL se ejercita en \rutarepo{crates/infrastructure/tests/postgres_document_queries.rs} y el contrato HTTP en \rutarepo{crates/web/tests/document_queries.rs}.\par}

\begin{table}[H]
  \centering
  \caption{Condiciones de comprobación de las consultas documentales por expediente.}
  \label{tab:pruebas-consultas}
  \begin{tabularx}{\linewidth}{@{}L{3.2cm} X@{}}
    \toprule
    \textbf{Frontera} & \textbf{Condiciones verificables} \\
    \midrule
    Entrada & Tamaño de página entre 1 y 100; nombre de hasta 200 caracteres sin controles; preservación literal de \texttt{\%} y \texttt{\_}; parámetros desconocidos rechazados. \\
    Autorización & Sesión vigente; Cliente denegado; pertenencia, rol y estado activo recargados; documento ajeno e inexistente con el mismo error. \\
    Persistencia & Filtros aplicados antes de paginar; orden por UUID; \texttt{has\_more} coherente; consultas que no decodifican vault ni evidencia. \\
    Auditoría & Un evento de listado o detalle por consulta aceptada; un fallo de auditoría impide devolver metadatos y conserva la cadena. \\
    HTTP & DTO sin criptograma ni material criptográfico, transferencia exacta de filtros y aislamiento de la asociación expediente/documento. \\
    \bottomrule
  \end{tabularx}
\end{table}

Las pruebas del navegador Qadra se conservan bajo \rutarepo{web/apps/qadra/tests/}. Las pruebas con respuestas HTTP simuladas permiten aislar formularios, navegación, permisos visibles, estados de error y solicitudes. El 14 de septiembre, \rutarepo{scripts/web-demo.sh} aprobó un escenario de navegador contra la API, PostgreSQL, Redis y OpenSSL reales, definido en \rutarepo{web/apps/qadra/tests/live/document-workflow.spec.mjs}. El recorrido inició sesión con contraseña y recuperación de segundo factor, creó un expediente desde la interfaz, cargó un documento, obtuvo sus metadatos, confirmó su sello, verificó su integridad y descargó un ZIP cuyos bytes se contrastaron con el archivo original. Tras cerrar sesión, recargar y autenticar de nuevo, recuperó el expediente y su documento desde el repositorio persistido. El escenario también comprobó ausencia de errores de JavaScript y de desbordamiento horizontal a 390 píxeles de ancho. Esta ejecución demuestra ese recorrido concreto; no sustituye los escenarios restantes de los 17 casos de uso ni la evaluación de usabilidad con usuarios del despacho.

\subsection{Versiones documentales inmutables}\label{sec:pruebas-versiones}

Las pruebas de versiones separan la preservación de una instantánea de la elección de la versión actual. La \cref{tab:pruebas-versiones} organiza sus escenarios: éxito, rechazo, concurrencia y recuperación. Las pruebas de aplicación verifican el orden de autorización, preparación y confirmación mediante puertos; las de infraestructura utilizan PostgreSQL y criptografía reales. La trazabilidad por archivo se presenta en la \cref{sec:matriz-versiones}, separada de la de consultas (\cref{sec:matriz-consultas}).

\begin{table}[H]
  \centering
  \caption{Escenarios de comprobación de versiones documentales.}
  \label{tab:pruebas-versiones}
  \begin{tabularx}{\linewidth}{@{}L{3.0cm} X@{}}
    \toprule
    \textbf{Propiedad} & \textbf{Escenarios positivos y negativos} \\
    \midrule
    Contexto cifrado & Versiones con contenido igual y distinto se descifran bajo su UUID y número; solicitan DEK propias. Sustituir vault, UUID, versión o digest produce rechazo. \\
    Concurrencia & Cuatro anexados con la misma cabeza esperada producen un ganador y conflictos; sellar la versión anterior conserva su evidencia. La inserción espera al predecesor y rechaza el hueco si este revierte. \\
    Atomicidad & Fallos de inserción, auditoría y commit diferido preservan documento y cadena anteriores. Fallar el evento de historial impide entregar sus metadatos. \\
    Consultas y acceso & Filtros sobre la cabeza, historial descendente, cursor exclusivo, primera versión disponible y errores de versión ausente; Cliente denegado y pertenencia o rol revocados. \\
    Continuidad & Migración repetida conserva filas y evidencia; origen con versión 7 admite la 8 y reconcilia la 7. Esquema incompleto, huecos y privilegios capaces de alterar la raíz se rechazan al abrir. \\
    Frontera HTTP & Número explícito en acciones y respuestas; alias ambiguos rechazados, cuerpo acotado, parámetros inválidos y conflictos tipados sin redirigir a otra versión. \\
    \bottomrule
  \end{tabularx}
\end{table}

La demostración \rutarepo{scripts/api-version-demo.sh}, integrada en \rutarepo{scripts/api-demo.sh}, añade una segunda versión a un documento importado, comprueba el conflicto con la cabeza anterior y recorre historial, sellado, verificación y exportación explícitos. Compara el ZIP histórico con el obtenido antes del anexado y, tras restaurar PostgreSQL, compara nuevamente las exportaciones. OpenSSL valida firma y sello del contenido nuevo. Estas comprobaciones verifican las muestras y restauraciones ensayadas; una copia antigua íntegra aún puede conservar validez interna si no existe anclaje externo de la cabeza de auditoría.

La validación final de Qadra aprobó 28 pruebas unitarias y 32 de navegador con respuestas simuladas, además de formato y compilación, bajo Node~22.22.2 y npm~10.9.7. El escenario real de \rutarepo{scripts/web-demo.sh} también aprobó: selló la versión 1 y descargó su ZIP, añadió una versión 2 con nombre y contenido distintos, la selló y verificó, y contrastó la descarga con los bytes del archivo nuevo. Al seleccionar nuevamente la versión 1 obtuvo un ZIP idéntico al inicial. Después de cerrar sesión, recargar y autenticar de nuevo, recuperó ambas versiones. No se registraron errores JavaScript. Las pruebas simuladas aíslan además conflictos de carga y respuestas tardías de sesiones o selecciones anteriores. Estos resultados se contabilizan por separado de Rust y no sustituyen la evaluación de usabilidad con personal del despacho.

\subsection{Clasificación y carga documental atómica}\label{sec:pruebas-clasificacion}

La clasificación mantiene una historia organizativa independiente de las versiones cifradas. Sus pruebas distinguen validación del valor, orquestación de permisos, persistencia transaccional, adaptación HTTP y comportamiento del navegador. La \cref{tab:pruebas-clasificacion} organiza las condiciones definidas por los archivos de prueba; la trazabilidad se detalla en la \cref{sec:matriz-clasificacion}. El canon de metadatos es un formato propio del prototipo, por lo que sus vectores fijan interoperabilidad entre Rust, PostgreSQL y OpenSSL y no se presentan como vectores oficiales de un estándar externo.

\begin{table}[H]
  \centering
  \caption{Condiciones de comprobación de clasificación documental.}
  \label{tab:pruebas-clasificacion}
  \begin{tabularx}{\linewidth}{@{}L{3.2cm} X@{}}
    \toprule
    \textbf{Frontera} & \textbf{Condiciones verificables} \\
    \midrule
    Valores y canon & Límites por caracteres Unicode, controles antes de recortar, conteo previo a deduplicar, etiquetas exactas y vectores con longitudes UTF-8. \\
    Autorización & Cuenta y pertenencia vigentes, Cliente denegado, UUID ajeno indistinguible del ausente y ambos permisos de carga clasificada revalidados antes del commit. \\
    Transacción & Un sucesor por revisión esperada; carga inicial y dos eventos indivisibles; fallo de fila, segundo evento o commit diferido sin residuos. \\
    Concurrencia & Clasificaciones competidoras producen un ganador; anexado de contenido conserva la clasificación de su commit; revocación confirmada antes de autorizar impide escribir. \\
    Consultas & Cabezas de archivo y clasificación seleccionadas antes de filtrar; retirada de etiquetas antiguas, coincidencia literal y lectura sin decodificar evidencia. \\
    Persistencia & Canon y digest Rust/SQL coincidentes, autor UUID y correo histórico, UTC con nanosegundos, rechazo de datos almacenados inconsistentes y privilegios restrictivos. \\
    HTTP & JSON estricto, límites de parte y petición, orden libre, partes duplicadas o incompletas rechazadas y fallo tardío del cuerpo sin carga confirmada. \\
    \bottomrule
  \end{tabularx}
\end{table}

La prueba de transporte entrega el cierre multipart antes de un epílogo fragmentado o de un fallo posterior del flujo. La condición exigida es que el límite se aplique a todos los bytes recibidos y que el caso de uso no comience antes de completar satisfactoriamente el cuerpo. Las pruebas de clasificación también separan la revisión cero de una decisión posterior con valores vacíos y exigen conservar el correo capturado aunque la cuenta cambie de dirección. La integridad de esas filas no se equipara con una firma individual ni con detección universal de una restauración administrativa antigua.

La validación final de la interfaz para clasificación aprobó 34 pruebas unitarias, 47 escenarios con respuestas HTTP simuladas y dos recorridos con servicios reales, además de formato y compilación. El recorrido de \rutarepo{web/apps/qadra/tests/live/document-classification.spec.mjs} carga archivo y revisión uno, clasifica de nuevo y genera un conflicto desde dos contextos autenticados. El segundo formulario conserva su borrador, compara la revisión concurrente y confirma una nueva decisión. Añadir una versión de contenido mantiene la clasificación actual; vaciar todos los campos crea la revisión cinco y conserva las cinco decisiones históricas con autor. Los filtros dejan de encontrar etiquetas retiradas y el reingreso recupera las versiones y revisiones. Los ZIP de las versiones 1 y 2 permanecen idénticos después de reclasificar y vaciar valores. El recorrido de versiones también aprobó, sin errores JavaScript. Estas comprobaciones de navegador se cuentan por separado de Rust y no sustituyen la evaluación de usabilidad.

{\emergencystretch=3em La restauración inicial reveló que las funciones SQL de canon dependían del \texttt{search\_path} de la sesión; la resolución explícita de sus referencias contra el esquema de instalación (\cref{sec:impl-http-migracion}) elimina esa dependencia. Las regresiones de \rutarepo{crates/infrastructure/tests/document_metadata_restore.rs} ejecutan las funciones con ruta de búsqueda vacía y un ciclo real de \texttt{pg\_dump}/\texttt{pg\_restore}; comprueban conservación de filas, reapertura con rol restringido, anexado posterior y rechazo de una huella incorrecta. La demostración HTTP confirma también el inventario completo y los ZIP anteriores después de restaurar.\par}

\subsection{Validación y aceptación de la admisión documental general}\label{sec:pruebas-admision-documental}

Las comprobaciones focales de la admisión general
(\cref{sec:impl-admision-documental}) distinguen la política de aplicación, el rechazo HTTP y la ejecución de analizadores nativos. La
\cref{tab:pruebas-admision-documental} presenta los resultados aprobados por
frontera; no representa una nueva ejecución global del proyecto ni actualiza
las mediciones históricas de cobertura. Los 29 casos de infraestructura y los
cinco del adaptador nativo son pruebas distintas, sin acumular sus repeticiones
focales como casos adicionales.

\begin{table}[H]
  \centering
  \caption{Pruebas focales de admisión de nuevas cargas y versiones.}
  \label{tab:pruebas-admision-documental}
  \begin{tabularx}{\linewidth}{@{}L{3.2cm} r X@{}}
    \toprule
    \textbf{Frontera} & \textbf{Aprobadas} & \textbf{Comprobación principal} \\
    \midrule
    Aplicación & 8 & Tres ingresos documentales, rechazo previo a persistencia y reautenticación. \\
    Cliente Node & 12 & Errores tipados, contexto de sesión y envío único. \\
    Navegador controlado & 8 & Rechazo de carga o versión y conservación de archivo y borrador. \\
    Adaptador nativo & 5 & Familias reales y corrupción interna de contenido comprimido. \\
    Proceso privado & 5 & Límites de recursos, descriptores heredados y protocolo. \\
    Infraestructura & 29 & Estructuras de archivo, inventario multimedia y supervisión de procesos. \\
    HTTP & $1+1$ & Una prueba unitaria y una de rutas: categorías públicas y transporte de errores. \\
    Instalador & 15 & Fuente, configuración y provisión comprobadas en VPS3. \\
    \bottomrule
  \end{tabularx}
\end{table}

Las pruebas nativas comprueban que no basta una envoltura válida cuando el
contenido comprimido interno está dañado. Las del proceso privado observan sus
límites, la retirada de descriptores ajenos y el rechazo de respuestas de
protocolo incompletas o inesperadas. El supervisor diferencia exceso de salida,
vencimiento, rechazo del archivo e indisponibilidad de ejecución. Los recorridos
con HTTP controlado comprueban que el fallo no borra el archivo seleccionado ni
los metadatos y que no provoca una segunda escritura automática. Esos recorridos
no sustituyen el ingreso contra PostgreSQL, Redis y el decodificador compuesto
en el servidor real.

La aceptación integrada mediante \rutarepo{scripts/api-demo.sh} terminó con
salida cero. Después del respaldo y la restauración, las ocho familias
conservaron sus bytes originales exactos. El navegador con servicios reales
aprobó cuatro escenarios en 27.2~s: uno nuevo de admisión documental, cuya
ejecución duró 8.4~s, y tres regresiones de contenido. Se inspeccionaron las
capturas de escritorio a 1440 píxeles y móvil a 390 píxeles, sin desbordamiento
horizontal. Estas evidencias complementan las focales; no acreditan una
campaña de usabilidad con personas.

La dependencia se verificó con el usuario del runner en los tres VPS. VPS1,
con AlmaLinux, requirió una compilación nativa porque el binario preparado en
Ubuntu~22 depende de GLIBC~2.35. La comprobación posterior de cada instalación
es de solo lectura. Clippy aprobó el workspace y todos los targets con
advertencias denegadas. Los gates globales y la integración no forman parte de
esta evidencia local: la integración del incremento en \texttt{main}, con su
confirmación, consta en \rutarepo{docs/product-completion.md}, y las cifras
anteriores describen la aceptación original y no una nueva ejecución. Los
resultados no acreditan detección de malware, autenticidad jurídica ni
ampliación de la política de soportes PDF/DOCX (\cref{sec:impl-formatos}). El contrato y la
distinción entre evidencia local y cierre global se conservan en
\rutarepo{docs/document-upload-admission-api.md} y en
\rutarepo{docs/adr/0058-bounded-general-document-admission.md}.


\subsection{Contenido exacto y buzón interno de integridad}\label{sec:pruebas-contenido-documental}

Las pruebas de la entrega de contenido descrita en la \cref{sec:impl-contenido-documental} separan autorización, procesamiento criptográfico, commit auditado y transporte. Las primeras ejecuciones fallidas acreditaron la ausencia del comportamiento antes de implementarlo. La focal del 19 de septiembre de 2026 aprobó trece casos de aplicación y siete HTTP en 41.997~s, con fuentes estables de las capas ejercitadas. Comprueban entrega histórica pendiente de sello, rechazo de otro ámbito o versión, reautenticación del principal completo, cota de lectura, fallo de auditoría y ausencia de bytes cuando falla la validación o el registro del incidente. La capa HTTP distingue el rechazo de integridad de una indisponibilidad y verifica adjunto, cabeceras exactas y consultas acotadas.

El cliente aprobó diecinueve pruebas Node en 0.563~s: comprobación de contexto y cabeceras, invalidación de respuestas tardías, valores de incidente, instantes y cursor. El navegador con HTTP controlado aprobó trece de catorce escenarios en 58.274~s; la espera del caso de descarga durante cierre de sesión agotó su límite. Ese comando no se presenta como aprobado. La repetición focal del único caso pendiente aprobó en 9.783~s tras corregir la espera de una respuesta descartada, con las fuentes estables: son catorce escenarios distintos aprobados en ejecuciones separadas. Los siete escenarios del buzón comprobaron escritorio y móvil, actualización después de error, logout y ausencia de acceso para los otros tres roles.

Ocho pruebas PostgreSQL aprobaron en 42.427~s con servicios desechables: contenido histórico, cota previa a materialización, permisos, rollback, idempotencia, catálogo, inventario y restauración. Dos casos HTTP adicionales reprodujeron una capacidad liberada antes de soltar los bytes; aprobaron junto con las cuatro regresiones de contenido en 24.454~s después de vincular el cupo al último propietario del buffer. Estas focales no constituyen una regresión global.

El navegador con servicios reales aprobó tres recorridos en 285.120~s de comando, de los cuales 24.6~s corresponden a Playwright, con 2490 fuentes estables. Los recorridos Owner a 1440 y 390 píxeles descargaron los bytes exactos de una versión V1 pendiente de sello después de registrar V2, comprobaron sus cabeceras y resumen, y consultaron el aviso, el buzón persistente y la apertura de la versión observada. El expediente del recorrido móvil estaba cerrado administrativamente. El tercer recorrido comprobó la lectura del Paralegal asignado, la denegación de Client y del buzón a otros roles, así como la eliminación de la captura visible después de revocar la asignación.

La preparación del ensayo alteró dos vaults AES aislados para producir rechazos reales y registrar sus incidentes; un bloque de finalización restauró los bytes originales antes de iniciar el navegador. Esta operación pertenece al fixture desechable y no constituye una reparación ofrecida por el producto. Los incidentes persistieron y las denegaciones por permisos no generaron observaciones nuevas. Se inspeccionaron cuatro capturas de contenido y buzón en escritorio y móvil y dos ampliaciones del detalle: los elementos de Qadra resultaron legibles, sin recortes ni desbordamiento. La muestra no constituye una evaluación de usabilidad con personas.

La cuarta campaña API completa aprobó en 372.048~s con 2490 fuentes estables, después de corregir la conexión SQL del fixture. Comprobó V1/V2 pendientes, bytes y cabeceras, cuatro roles, cierre, revocación y rechazo de validación con incidente atómico. La corrupción desechable se revirtió antes del respaldo. La restauración conservó los bytes exactos y las respuestas históricas del incidente, con 19 raíces documentales, 25 versiones y un incidente inmutable; el prefijo importado de cuatro documentos y 63 eventos retuvo un ZIP idéntico. Los tres intentos anteriores permanecen fallidos en el informe. El cierre de CI y la compilación e inspección del PDF conservan verificaciones independientes. Los resultados y las limitaciones se documentan en \rutarepo{docs/verification-report.md}. La descarga no evalúa la identidad del firmante; los avisos tampoco acreditan un ataque ni atención humana.


\subsection{Firma documental individual con clave externa}
\label{sec:pruebas-firma-individual}

La verificación local del 10 de octubre de 2026 comprobó el vínculo de una
cuenta Litigator con un certificado público y la firma de una versión PDF
exacta mediante dos pruebas externas. El registro canónico identifica la
cuenta y su certificado; el consentimiento documental incorpora expediente,
documento, versión, operación, confianza y resumen de la firma del PDF.
La identidad de la cuenta no se deduce del nombre común del certificado.
Los resultados se separan por frontera en la
\cref{tab:matriz-firma-individual}; no constituyen una nueva medición global
de cobertura ni reemplazan los cortes históricos de este capítulo.

\subsubsection{Valores, emisión y pruebas criptográficas}

Doce pruebas de dominio aprobaron los vectores literales de 142 bytes del
vínculo \texttt{DOCBIND1} y de 258 bytes del consentimiento \texttt{DOCSIGN1},
su reconstrucción estricta y el rechazo de identidades, contadores y propósitos
alterados. Cinco pruebas de emisión mediante la CLI aprobaron con OpenSSL:
la hoja conserva la clave pública del CSR y el perfil interno de firma,
sin generar una clave privada del titular dentro de la CA. Se rechazaron
solicitudes alteradas, excesivas o privadas, parámetros RSA no admitidos,
sobrescritura de archivos y sustitución de la clave pública emitida.

Ocho casos criptográficos comprobaron las dos firmas RSA de 3\,072 bits,
certificado, confianza exacta, vigencia y revocación. La firma del
consentimiento no sustituye la comprobación de la firma separada del PDF.
Otro documento, versión, cuenta o propósito no puede reutilizar esa intención.
Once casos del servicio de vinculación y catorce del servicio documental
comprobaron reautenticación del principal completo, recibos exactos, cambios
de cuenta o confianza, retiro, respuesta perdida y ausencia de una segunda
escritura en la repetición autorizada.

Ocho pruebas del procesador comprobaron la conservación del material personal,
la TSA configurada, el rechazo de sellos de tiempo inválidos y la evaluación
del certificado en el instante de aceptación. Seis casos de evidencia
histórica rechazaron cambios de material, contexto o formato. La exportación
aprobó once pruebas, incluidas las cuatro de la extensión individual; el
repositorio de archivos aprobó diez, conservando la lectura de sellos anteriores.
Siete pruebas HTTP con puertos controlados verificaron rutas, límites,
identificadores, errores y la distinción explícita de evidencia del servidor.
Estos ensayos controlados no acreditan por sí solos el recorrido de servicios
reales descrito a continuación.

\subsubsection{Atomicidad, concurrencia y restauración nativa}

Con PostgreSQL~18.6 aislado y roles operativos restringidos aprobaron diecisiete
casos distintos: siete de vínculos, cinco del registro documental y cinco de
invariantes. El conteo reúne la campaña inicial y las repeticiones focales de
los casos corregidos, sin contar dos veces una identidad de prueba. Se
comprobaron permisos actuales, cambios de cuenta o confianza, huellas retenidas
por su cuenta original, un único vínculo activo, retiro terminal, concurrencia
y rollback de documento, recibo y auditoría ante un fallo de escritura.

Una evidencia personal sin recibo y una divergencia entre ambos se rechazaron
al confirmar la transacción. Desactivar cualquiera de los dos disparadores
diferidos impidió reabrir el adaptador. El volcado y la restauración lógica
conservaron la evidencia y sus eventos después de retirar el vínculo. Una
regresión reprodujo además la sustitución de la CRL capturada por una sucesora
válida: el vínculo se había registrado con la primera revisión y el documento
firmado con la segunda. El inventario corregido exige que la confianza
embebida coincida con la revisión exacta del recibo, y rechazó la sustitución
por la tercera. La concurrencia nativa conservó un solo recibo. En aplicación
se comprobó además que prevalecía la primera evidencia aceptada aunque ambas
solicitudes prepararan sellos de tiempo distintos.

Los primeros intentos de migración y la prueba que suponía un nombre automático
de restricción no se contabilizan como aprobaciones. Sus correcciones se
verificaron antes del resultado anterior. La admisión de formato y la TSA de
estas pruebas nativas utilizan puertos controlados; su finalidad es comprobar
el contrato transaccional con PostgreSQL real.

\subsubsection{Aceptación HTTP, evidencia pública y continuidad}

El recorrido focal HTTP aprobado utilizó PostgreSQL~18.6, Valkey~8.1.10 como
almacén compatible con Redis, ingreso real por contraseña y MFA, CA interna,
OpenSSL y TSA local real. La clave sintética se generó en un directorio del
cliente separado de la CA; el servidor recibió certificado y pruebas públicas.
La campaña ejecutó los pasos de
\rutarepo{scripts/api-individual-signing-demo.sh} y comprobó independientemente
los bytes canónicos mediante su auxiliar Python.

El litigante vinculó su certificado, descargó el contenido exacto de la versión
uno y aportó sus dos firmas. La versión dos del mismo documento permaneció
sin sellar. Se descartó deliberadamente la respuesta de confirmación y se
recuperó el recibo por su operación; la repetición devolvió el mismo resultado
sin otro registro de firma. OpenSSL verificó las pruebas públicas del ZIP,
la cadena y CRL capturadas y el token RFC~3161 sobre el PDF.

Después se publicó una CRL que revocaba la hoja. La verificación mantuvo la
validez histórica y mostró por separado la revocación actual; una nueva firma
fue rechazada. El retiro del vínculo conservó la evidencia original. Tras
respaldar y restaurar PostgreSQL, invalidar las sesiones anteriores e ingresar
con MFA nueva, se recuperaron los mismos recibos, instantes y bytes del ZIP.
El documento sellado mediante el firmante del servidor mantuvo su clasificación
histórica y su paquete sin cambios.

El corte de interfaz aprobó 872 pruebas Node en 4.169 segundos y seis escenarios
de navegador con HTTP controlado en 19.4 segundos, con un trabajador. Incluyó
revisión del PDF exacto, recuperación explícita, restricciones por rol, retiro
móvil y renovación de la confirmación tras rechazar una firma. La compilación
del workspace y la compilación web aprobaron.

La aceptación posterior de Qadra con servicios reales aprobó dos recorridos
en 20.2 segundos de Playwright, con un trabajador: escritorio de 1440 por
1000 píxeles y móvil de 390 por 844. El navegador descargó el PDF y los bytes
canónicos, recibió las firmas generadas externamente con OpenSSL y confirmó
la operación. La prueba interceptó una respuesta exitosa después del commit
para simular su pérdida. Tras revocar la sesión, observar el rechazo de acceso
e ingresar con MFA nueva, la interfaz recuperó el recibo original mediante
consulta, sin repetir el envío. Descargó el ZIP y, después de retirar el
vínculo, mostró la validez histórica y el retiro posterior por separado.

La primera ejecución detectó una carrera del instrumento: leía la respuesta
capturada antes de terminar la intercepción HTTP. La corrección espera su
finalización conservando la exigencia de respuesta exitosa y recibo exacto.
Los dos recorridos finales aprobaron sin errores de página. Esta aceptación
de navegador no repite la revocación por CRL ni la restauración SQL del
recorrido HTTP anterior. Las ocho capturas se inspeccionaron para comprobar
legibilidad de controles, referencias y resultados en ambos tamaños. La
visualización embebida del PDF no se acredita mediante capturas headless;
los bytes descargados y sus firmas se comprobaron por separado. Los controles
de la revisión publicada y el despliegue conservan su evidencia propia.

La TSA local acredita el resumen del contenido bajo su confianza configurada;
no prueba que las firmas personales existieran al generar el token. Registro
y aceptación conservan los instantes de la aplicación. Las claves sintéticas
permiten comprobar el protocolo, sin demostrar custodia exclusiva de una
persona, identidad civil, FIREL ni una constancia NOM-151 emitida por un PSC.

